\chapter{Các công trình liên quan}
\label{chap:related-work}

Chương này khảo sát nền tảng lý thuyết và các công trình tiền đề mà khóa luận
dựa vào hoặc lấy cảm hứng. Chúng tôi đi từ những khái niệm nền tảng nhất (sổ cái
phân tán, hàm băm, cây Merkle) đến các hệ thống blockchain tiêu biểu (Bitcoin,
Ethereum, Hyperledger Fabric), rồi tới các thuật toán đồng thuận, công nghệ máy
ảo hợp đồng thông minh và hạ tầng mạng ngang hàng. Cuối chương, chúng tôi định
vị hệ thống của khóa luận trong bức tranh tổng thể đó.

\section{Sổ cái phân tán và hàm băm mật mã}
\label{sec:dlt-hash}

Cốt lõi của mọi blockchain là một sổ cái phân tán: một danh sách các
bản ghi (giao dịch) được nhóm thành các \textit{block} (khối), trong đó mỗi
block chứa ``dấu vân tay'' mật mã của block liền trước. Dấu vân tay này được
tính bằng một hàm băm mật mã (cryptographic hash function) --- một hàm
biến dữ liệu có độ dài tùy ý thành một chuỗi bit có độ dài cố định, với ba thuộc
tính then chốt: kháng tiền ảnh (preimage resistance), kháng tiền ảnh thứ hai, và
kháng va chạm (collision resistance). Một thay đổi dù chỉ một bit ở dữ liệu đầu
vào cũng làm thay đổi hoàn toàn giá trị băm, và về mặt tính toán là không khả
thi để đảo ngược.

Vì mỗi block tham chiếu tới hash của block trước, việc sửa đổi một block cũ sẽ
làm thay đổi hash của nó, kéo theo phải tính lại toàn bộ các block sau --- điều
khiến dữ liệu trở nên gần như bất biến. Đây chính là cơ chế ``chuỗi''
(chain) trong từ ``blockchain''. Hệ thống trong khóa luận sử dụng SHA-256 (cụ
thể là double-SHA256, giống Bitcoin)~\cite{fips180,nakamoto2008bitcoin} cho hash
block.

\subsection{Cây Merkle}

Để tóm tắt toàn bộ các giao dịch trong một block thành một giá trị duy nhất,
blockchain sử dụng cây Merkle~\cite{merkle1979} (Merkle tree, hay hash
tree). Trong cấu trúc này, các giao dịch là lá; mỗi nút trong là hash của hai
nút con ghép lại; và nút gốc --- \textit{Merkle root} --- tóm tắt toàn bộ cây.
Cây Merkle mang lại hai lợi ích: (1)~bất kỳ thay đổi nào ở một giao dịch đều làm
thay đổi Merkle root, nên root đóng vai trò bằng chứng toàn vẹn; và (2)~người ta
có thể chứng minh một giao dịch thuộc về block mà chỉ cần cung cấp $O(\log n)$
hash trên đường từ lá đến gốc, thay vì tải toàn bộ block --- cơ sở cho các
``light client'' (máy khách nhẹ). Hệ thống của khóa luận tính Merkle root từ
danh sách \texttt{txid} của các giao dịch và đưa vào header của block.

\section{Bitcoin và mô hình UTXO}
\label{sec:bitcoin-utxo}

Bitcoin~\cite{nakamoto2008bitcoin} là blockchain đầu tiên giải quyết được bài
toán ``chi tiêu hai lần'' (double-spending) trong môi trường không cần tin cậy,
thông qua cơ chế đồng thuận Proof-of-Work và một cấu trúc dữ liệu giao dịch độc
đáo: mô hình UTXO (Unspent Transaction Output --- đầu ra giao dịch chưa
tiêu).

Thay vì lưu ``số dư tài khoản'' như ngân hàng truyền thống, Bitcoin lưu một tập
hợp các ``tờ tiền'' chưa tiêu. Mỗi giao dịch gồm:

\begin{itemize}
  \item Các đầu vào (inputs, ký hiệu \texttt{Vin}): tham chiếu và
  ``tiêu'' các UTXO đã tồn tại từ trước.
  \item Các đầu ra (outputs, ký hiệu \texttt{Vout}): tạo ra các UTXO
  mới, gán cho người nhận.
\end{itemize}

Số dư của một địa chỉ chính là tổng giá trị các UTXO mà địa chỉ đó kiểm soát. Mô
hình UTXO có ưu điểm về khả năng song song hóa kiểm tra giao dịch và tính riêng
tư, nhưng kém trực quan hơn mô hình tài khoản (account model) khi biểu diễn
trạng thái phức tạp. Hệ thống của khóa luận hỗ trợ mô hình UTXO: world
state ở Committing Peer được tổ chức như một tập UTXO, và CLI client đóng vai trò
một ví UTXO đầy đủ.

\section{Ethereum, mô hình tài khoản và hợp đồng thông minh}
\label{sec:ethereum}

Ethereum~\cite{buterin2014ethereum,wood2014ethereum} mở rộng ý tưởng blockchain
từ một hệ thống thanh toán thành một ``máy tính thế giới'' (world computer) có
khả năng lập trình. Đóng góp lớn nhất của Ethereum là hợp đồng thông
minh (smart contract): các chương trình được lưu và thực thi trên blockchain,
tự động chạy khi có giao dịch gọi tới.

Ethereum sử dụng mô hình tài khoản (account-based model) thay cho UTXO:
trạng thái toàn cục là một ánh xạ từ địa chỉ tài khoản đến số dư và bộ nhớ hợp
đồng. Hợp đồng được biên dịch sang bytecode của \textit{Máy ảo Ethereum} (EVM)
và thực thi theo mô hình Order-Execute: các giao dịch được sắp xếp
trước (qua đồng thuận), sau đó \textit{mọi} node thực thi lại tuần tự để cập nhật
trạng thái. Mô hình này đơn giản về mặt khái niệm nhưng hạn chế về hiệu năng và
khả năng song song: mọi node phải chạy lại mọi hợp đồng, và một hợp đồng độc hại
(ví dụ vòng lặp vô hạn) có thể gây nguy hiểm --- vấn đề mà Ethereum giải quyết
bằng cơ chế ``gas'' tính phí cho từng bước tính toán.

Hệ thống của khóa luận hỗ trợ hợp đồng thông minh (qua WebAssembly thay vì EVM)
nhưng áp dụng mô hình thực thi khác --- Execute-Order-Validate của Fabric --- sẽ
được phân tích ở phần sau.

\section{Hyperledger Fabric và mô hình Execute--Order--Validate}
\label{sec:fabric}

Hyperledger Fabric~\cite{androulaki2018hyperledger} là nền tảng blockchain có
cấp phép tiêu biểu cho doanh nghiệp, và là nguồn cảm hứng kiến trúc
chính của khóa luận. Đóng góp quan trọng nhất của Fabric --- so với mô hình
Order-Execute của Ethereum --- là kiến trúc Execute--Order--Validate,
tách quá trình xử lý một giao dịch thành ba pha độc lập:

\begin{enumerate}
  \item \textbf{Execute (Thực thi \& Xác nhận).} Giao dịch được gửi tới các
  \textit{endorsing peer}; mỗi peer thực thi hợp đồng (\textit{chaincode})
  trong môi trường cô lập và \textit{ký xác nhận} (endorse) kết quả, nhưng
  chưa ghi vào sổ cái. Vì các peer thực thi độc lập, bước này có thể
  được song song hóa.
  \item \textbf{Order (Sắp xếp).} Các giao dịch đã có chữ ký xác nhận được gửi
  tới \textit{ordering service}, nơi quyết định thứ tự toàn phần và gom chúng
  thành block. Đáng chú ý, ordering service không cần hiểu nội dung
  giao dịch --- nó chỉ sắp xếp, nên rất nhẹ và dễ mở rộng.
  \item \textbf{Validate (Kiểm tra \& Ghi).} Các \textit{committing peer} nhận
  block đã sắp xếp, kiểm tra chữ ký xác nhận có thỏa mãn chính sách endorsement
  không, kiểm tra xung đột phiên bản (đọc-ghi), rồi mới ghi vĩnh viễn.
\end{enumerate}

Mô hình này mang lại ba lợi ích cốt lõi: (a)~khả năng song song ở pha
thực thi; (b)~tính module hóa, cho phép thay thế cơ chế đồng thuận mà
không ảnh hưởng tới logic nghiệp vụ; và (c)~bảo mật theo lớp, vì mỗi
tầng chỉ thấy thông tin nó cần. Fabric thật sử dụng thư viện
etcd/raft cho ordering service.

Hệ thống của khóa luận tái hiện trung thành triết lý này: Core Service
đóng vai trò endorsing peer/gateway (pha Execute), Ordering Service đóng vai trò
ordering service (pha Order), và Committing Peer đóng vai trò committing peer
(pha Validate). Điểm khác biệt học thuật quan trọng là chúng tôi
tự cài đặt thuật toán Raft từ đầu thay vì dùng thư viện, và đơn giản
hóa một số khía cạnh (ví dụ chính sách endorsement) để tập trung vào việc làm
chủ cơ chế lõi. Bảng~\ref{tab:fabric-mapping} đối chiếu các vai trò.

\begin{table}[H]
\centering
\caption{Đối chiếu vai trò giữa Hyperledger Fabric và hệ thống của khóa luận}
\label{tab:fabric-mapping}
\begin{tabular}{@{}p{4cm}p{4.5cm}p{3.5cm}@{}}
\toprule
\textbf{Vai trò trong Fabric} & \textbf{Thành phần trong dự án} & \textbf{Pha} \\
\midrule
Endorsing Peer / Gateway & Core Service & Execute \\
Ordering Service (etcd/raft) & Ordering Service (Raft tự viết) & Order \\
Committing Peer & Committing Peer & Validate \\
Chaincode (Go/Java/Node) & Hợp đồng WASM (TinyGo) & Execute \\
\bottomrule
\end{tabular}
\end{table}

\section{Bài toán đồng thuận và các thuật toán tiêu biểu}
\label{sec:consensus}

Đồng thuận (consensus) là bài toán trung tâm của hệ phân tán: làm thế nào để
nhiều tiến trình, một số có thể bị lỗi, cùng thống nhất về một giá trị (hoặc một
chuỗi giá trị) duy nhất. Kết quả bất khả thi FLP~\cite{fischer1985flp} chứng minh
rằng trong một hệ bất đồng bộ hoàn toàn, không tồn tại thuật toán đồng thuận
\textit{tất định} vừa luôn an toàn vừa luôn kết thúc khi có dù chỉ một tiến
trình lỗi. Các thuật toán thực tế ``né'' kết quả này bằng cách dựa vào giả định
thời gian (timeout) để đảm bảo tính kết thúc (liveness) trong khi luôn giữ tính
an toàn (safety).

\subsection{Paxos}

Paxos~\cite{lamport1998paxos} của Leslie Lamport là thuật toán đồng thuận tất
định đầu tiên được chứng minh đúng đắn, và là nền tảng lý thuyết cho nhiều hệ
thống. Tuy nhiên, Paxos nổi tiếng là \textit{khó hiểu} và khó hiện thực đúng,
đặc biệt khi mở rộng thành Multi-Paxos để đồng thuận một chuỗi giá trị.

\subsection{Raft}
\label{sec:raft-related}

Raft~\cite{ongaro2014raft} được Ongaro và Ousterhout thiết kế với mục tiêu
tường minh là \textit{dễ hiểu} (understandability) mà vẫn tương đương Paxos về
năng lực. Raft phân rã bài toán đồng thuận thành ba bài toán con tách biệt:

\begin{itemize}
  \item \textbf{Bầu lãnh đạo (Leader Election):} tại mỗi thời điểm, chọn một
  node làm \textit{Leader}; các node khác là \textit{Follower}. Thời gian được
  chia thành các \textit{term} (nhiệm kỳ), mỗi term có tối đa một Leader. Raft
  chuẩn dùng timeout ngẫu nhiên: Follower nào hết giờ chờ trước sẽ trở
  thành \textit{Candidate} và xin phiếu bầu.
  \item \textbf{Nhân bản log (Log Replication):} chỉ Leader nhận lệnh từ client,
  ghi vào log của mình, rồi sao chép cho Follower. Khi một entry được sao chép
  tới đa số (majority) node, nó được coi là đã \textit{commit}.
  \item \textbf{An toàn (Safety):} một loạt ràng buộc (chỉ bầu cho candidate có
  log đủ mới, chỉ commit entry của term hiện tại...) đảm bảo không có hai Leader
  hợp lệ cùng term và log không bị rẽ nhánh.
\end{itemize}

Khái niệm quorum đa số ($\lfloor N/2 \rfloor + 1$ trên tổng $N$ node)
là then chốt: vì hai tập đa số bất kỳ luôn giao nhau ở ít nhất một node, không
thể có hai quyết định mâu thuẫn được commit. Nhờ vậy Raft chịu được tối đa
$\lfloor (N-1)/2 \rfloor$ node lỗi dừng. Raft là loại đồng thuận
chịu lỗi dừng (Crash Fault Tolerant --- CFT), giả định các node không
có hành vi độc hại, chỉ có thể dừng hoặc chậm.

Khóa luận tự cài đặt Raft nhưng thay cơ chế bầu cử bằng timeout ngẫu
nhiên bằng một biến thể bầu lãnh đạo theo độ ưu tiên (xem
Chương~\ref{chap:methodology}), nhằm loại bỏ hiện tượng chia phiếu và đạt hội tụ
xác định.

\subsection{PBFT và đồng thuận chịu lỗi Byzantine}

Trong môi trường mà các node có thể có hành vi \textit{tùy ý} (kể cả độc hại) ---
bài toán Tướng Byzantine~\cite{lamport1982byzantine} --- người ta cần các thuật
toán chịu lỗi Byzantine (Byzantine Fault Tolerant --- BFT).
PBFT~\cite{castro1999pbft} của Castro và Liskov là thuật toán BFT thực tế đầu
tiên, chịu được tối đa $\lfloor (N-1)/3 \rfloor$ node độc hại với chi phí truyền
thông $O(N^2)$ thông điệp mỗi vòng. So với CFT, BFT an toàn hơn trong môi trường
ít tin cậy nhưng đắt hơn nhiều về truyền thông và phức tạp hơn khi hiện thực.

Cachin và Vukoli\'c~\cite{cachin2017blockchain} đã khảo sát một cách hệ thống
các giao thức đồng thuận trong bối cảnh blockchain, chỉ ra rằng việc lựa chọn
CFT hay BFT phụ thuộc vào mô hình tin cậy của ứng dụng. Với một blockchain có
cấp phép trong nội bộ một liên minh doanh nghiệp đã định danh và chịu ràng buộc
pháp lý, mô hình CFT (như Raft) thường là đủ và được ưa chuộng vì hiệu năng cao
--- đây cũng là lựa chọn của Fabric (với Raft) và của khóa luận này.

\begin{table}[H]
\centering
\caption{So sánh các thuật toán đồng thuận}
\label{tab:consensus-compare}
\begin{tabular}{@{}p{2.4cm}p{2.2cm}p{2.6cm}p{3.4cm}@{}}
\toprule
\textbf{Thuật toán} & \textbf{Mô hình lỗi} & \textbf{Số node lỗi chịu được} & \textbf{Ghi chú} \\
\midrule
Paxos & Crash & $\lfloor(N-1)/2\rfloor$ & Khó hiểu, nền tảng lý thuyết \\
Raft & Crash & $\lfloor(N-1)/2\rfloor$ & Dễ hiểu; dùng trong dự án \\
PBFT & Byzantine & $\lfloor(N-1)/3\rfloor$ & Chi phí $O(N^2)$ \\
PoW (Bitcoin) & Byzantine & $<50\%$ sức tính & Tốn năng lượng \\
\bottomrule
\end{tabular}
\end{table}

\section{WebAssembly cho hợp đồng thông minh}
\label{sec:wasm-related}

Việc lựa chọn công nghệ thực thi hợp đồng thông minh ảnh hưởng trực tiếp tới ba
yêu cầu: tính xác định, tính cô lập và khả năng đa ngôn ngữ.
WebAssembly (WASM)~\cite{haas2017wasm,wasmspec} là một định dạng mã máy
ảo dạng stack, ban đầu được thiết kế cho trình duyệt web nhưng nhanh chóng trở
thành lựa chọn phổ biến cho các môi trường thực thi phía máy chủ và blockchain.
WASM thỏa cả ba yêu cầu: nó chạy trong một bộ nhớ tuyến tính cô lập (sandbox), có
ngữ nghĩa thực thi xác định, và có thể được biên dịch từ nhiều ngôn ngữ bậc cao
(Rust, Go, C/C++, AssemblyScript...).

Nhiều nền tảng blockchain hiện đại đã chuyển sang WASM: Polkadot dùng WASM cho
runtime, CosmWasm đưa hợp đồng WASM vào hệ sinh thái Cosmos, và Ethereum từng
nghiên cứu eWASM như một thay thế cho EVM. Để giao tiếp với hệ điều hành chủ,
WASM phía máy chủ thường dùng WASI (WebAssembly System
Interface)~\cite{wasi} --- một tập giao diện hệ thống chuẩn hóa.

Hệ thống của khóa luận biên dịch hợp đồng viết bằng Go sang WASM bằng
TinyGo, và thực thi bằng wazero --- một runtime WASM thuần Go
không phụ thuộc thư viện C, giúp đơn giản hóa việc build và triển khai. Giữa host
và contract có một giao diện hàm (ABI) quy ước để cấp phát bộ nhớ, truyền payload
và ghi world state, sẽ được mô tả chi tiết ở Chương~\ref{chap:methodology}.

\section{Mạng ngang hàng và lưu trữ}
\label{sec:p2p-storage}

\subsection{libp2p}

Các thành phần của một blockchain phân tán cần giao tiếp an toàn và linh hoạt.
Thay vì dùng HTTP/gRPC thông thường, hệ thống của khóa luận dùng
libp2p~\cite{libp2p} --- ngăn xếp mạng module hóa được tách ra từ dự án
IPFS~\cite{benet2014ipfs} và cũng là nền mạng của Ethereum 2.0. libp2p cung cấp:
(a)~định danh mật mã (PeerID suy ra từ khóa công khai); (b)~địa
chỉ tự mô tả (multiaddr) như \texttt{/ip4/127.0.0.1/tcp/6000/p2p/12D3Koo...};
(c)~ghép kênh (stream multiplexing) cho phép nhiều ``làn'' dữ liệu song
song trên một kết nối; và (d)~thương lượng giao thức qua các Protocol
ID. Những đặc tính này đặc biệt phù hợp với giao tiếp consensus, nơi cần gửi
nhiều loại thông điệp (heartbeat, đề xuất block, ACK...) một cách hiệu quả.

\subsection{LevelDB và lưu trữ nhúng}

Để lưu world state cần một cơ sở dữ liệu vừa nhanh vừa nhẹ.
LevelDB~\cite{leveldb} là một thư viện lưu trữ key-value nhúng, chạy
ngay trong tiến trình (không cần server riêng), dựa trên cấu trúc LSM-tree
(log-structured merge-tree) cho phép ghi tuần tự nhanh. Quan trọng nhất với
blockchain, LevelDB hỗ trợ ghi theo lô nguyên tử (atomic batch): nhiều
thao tác ghi/xóa hoặc được áp dụng toàn bộ, hoặc không gì cả --- thuộc tính then
chốt để cập nhật world state UTXO một cách nhất quán theo từng block.

\section{Đánh giá hiệu năng blockchain}
\label{sec:bench-related}

Việc đánh giá hiệu năng một blockchain có cấp phép đòi hỏi phương pháp luận cẩn
thận. BLOCKBENCH~\cite{dinh2017blockbench} là một trong những khung đánh giá đầu
tiên cho private blockchain, đề xuất đo throughput và độ trễ ở nhiều lớp khác
nhau của hệ thống. Một bài học quan trọng là cần phân biệt rõ \textit{nơi} đo:
độ trễ mà client cảm nhận (HTTP) khác với độ trễ ghi sổ thực sự (end-to-end).

Khi tốc độ gửi giao dịch vượt quá tốc độ xử lý bền vững của hệ thống, hàng đợi
giao dịch tích lũy thành backlog và độ trễ tăng vọt --- hiện tượng có
thể giải thích bằng định luật Little~\cite{little1961law}: số lượng
công việc trung bình trong hệ thống bằng tốc độ đến nhân với thời gian lưu trú
trung bình ($L = \lambda W$). Khóa luận áp dụng cách đo ba lớp (k6 đo HTTP,
submit đo lúc chấp nhận, commit đo lúc ghi sổ) và dùng định luật Little để lý
giải hành vi backlog ở điểm bão hòa.

\section{Định vị của khóa luận}
\label{sec:positioning}

Từ khảo sát trên, có thể định vị hệ thống của khóa luận như sau:

\begin{itemize}
  \item \textbf{Về kiến trúc:} tái hiện mô hình Execute--Order--Validate của
  Hyperledger Fabric~\cite{androulaki2018hyperledger}, nhưng đơn giản hóa để tập
  trung vào việc làm chủ cơ chế lõi.
  \item \textbf{Về đồng thuận:} tự cài đặt Raft~\cite{ongaro2014raft} từ đầu (so
  với Fabric dùng thư viện etcd/raft), với biến thể bầu lãnh đạo theo độ ưu tiên
  --- một đóng góp mang tính học thuật, đánh đổi tính tổng quát lấy tính đơn
  giản và hội tụ xác định.
  \item \textbf{Về hợp đồng thông minh:} dùng WebAssembly~\cite{haas2017wasm}
  (giống Polkadot/CosmWasm) thay vì EVM, với runtime thuần Go.
  \item \textbf{Về mô hình giao dịch:} hỗ trợ cả UTXO (giống
  Bitcoin~\cite{nakamoto2008bitcoin}) lẫn hợp đồng thông minh ghi world state
  (giống Ethereum~\cite{buterin2014ethereum}).
\end{itemize}

Khác biệt cốt lõi so với việc sử dụng lại một nền tảng có sẵn nằm ở
\textit{mục tiêu}: khóa luận không hướng tới một sản phẩm thương mại sẵn sàng
vận hành, mà hướng tới việc làm chủ tri thức nền tảng thông qua việc tự
hiện thực từng lớp. Chính vì vậy, báo cáo dành phần lớn dung lượng cho việc lý
giải các quyết định thiết kế và các điểm đánh đổi, đặc biệt ở tầng đồng thuận ---
nội dung được trình bày chi tiết trong Chương~\ref{chap:methodology}.
